% =============================================================================
%                    WEEK 9: MECHANISM DESIGN
% =============================================================================
\section{Week 9: Mechanism Design}

% -----------------------------------------------------------------------------
%                              9.1 TOPIC SUMMARY
% -----------------------------------------------------------------------------
\subsection{Topic Summary}

\begin{keybox}[Learning Objectives]
By the end of this week, you will be able to:
\begin{enumerate}
  \item Describe core concepts of \textbf{mechanism design}: social choice functions, incentive compatibility, and strategyproofness.
  \item Define the four classic \textbf{single-good auction formats} (English, Dutch, first-price sealed-bid, Vickrey) and their strategic equivalences.
  \item Model auctions as \textbf{Bayesian games} and compute Bayes-Nash equilibrium bid functions.
  \item Prove that truthful bidding is a \textbf{weakly dominant strategy} in the Vickrey (second-price) auction.
  \item State and apply the \textbf{Revenue Equivalence Theorem}.
  \item Analyse \textbf{collusion}, \textbf{cheating auctioneers}, and the \textbf{winner's curse}.
  \item Define the \textbf{VCG mechanism} for combinatorial auctions: allocation rule, transfer rule, and key properties.
  \item Describe \textbf{multi-unit auctions}, \textbf{combinatorial auctions}, and \textbf{continuous double auctions}.
\end{enumerate}
\end{keybox}

\separator

\subsubsection*{The Big Picture}

Weeks~7--8 studied fair division, where a planner distributes resources without payments. This week moves to \textbf{mechanism design}: engineering rules so that self-interested agents, each maximising their own utility, produce \textbf{desirable outcomes} (efficiency, revenue, truthfulness). The primary application is \textbf{auctions}---protocols for allocating goods among strategic bidders with private valuations. The crown jewel is the \textbf{VCG mechanism}, which achieves efficiency and truthfulness in general combinatorial settings.

\separator

\subsubsection*{What Examiners Like to Ask}

\begin{itemize}
  \item Fill in the payoff matrix for a \textbf{first-price} or \textbf{second-price} auction with given valuations and find Nash equilibria.
  \item Prove that \textbf{truthful bidding is weakly dominant} in the Vickrey auction.
  \item Compute the \textbf{BNE bid function} $s_i(\theta_i) = \frac{n-1}{n}\theta_i$ for the first-price auction with uniform priors.
  \item State and explain the \textbf{Revenue Equivalence Theorem}.
  \item Compute \textbf{VCG allocations and payments} for multi-item or multi-unit settings.
  \item Analyse \textbf{collusion}: optimal bids and surplus for colluding bidders in a second-price auction.
  \item Determine the \textbf{equilibrium price} in a continuous double auction.
  \item Check whether a VCG outcome is \textbf{individually rational} and \textbf{weakly budget-balanced}.
\end{itemize}

\separator

% -----------------------------------------------------------------------------
%           9.2 OVERVIEW + CORE CONCEPTS + EXTENDED THEORY
% -----------------------------------------------------------------------------
\subsection{Overview + Core Concepts + Extended Theory}

\subsubsection*{Roadmap}
\begin{enumerate}
  \item Mechanism design: motivation and framework
  \item Single-good auction formats (English, Dutch, first-price, Vickrey)
  \item Bayesian games and Bayes-Nash equilibrium
  \item First-price auction: BNE analysis
  \item Vickrey auction: strategyproofness
  \item Strategic equivalences and Revenue Equivalence Theorem
  \item Collusion, cheating auctioneers, winner's curse
  \item Multi-unit and combinatorial auctions
  \item VCG mechanism
  \item Continuous double auctions
\end{enumerate}

\bigseparator

% --- Mechanism design intro ---
\subsubsection{Mechanism Design: Motivation}

\begin{defbox}[Mechanism Design]
\textbf{Mechanism design} is the ``inverse'' of game theory: instead of analysing a given game, we \textbf{design the rules} (action spaces, information flows, payoff rules) so that certain objectives are achieved at equilibrium when agents maximise their own utility.

Objectives include: \textbf{efficiency} (maximise total welfare), \textbf{revenue maximisation}, \textbf{truthfulness} (incentive compatibility), \textbf{fairness}.
\end{defbox}

Examples:
\begin{itemize}
  \item \textbf{Elections:} design voting protocols so that the best candidate wins or dissatisfaction is minimised.
  \item \textbf{Marketplaces:} design auctions to maximise efficiency, seller profit, or truthful reporting.
\end{itemize}

\separator

% --- Auctions ---
\subsubsection{Single-Good Auctions}

\begin{defbox}[Auction Setting]
One indivisible item, one seller, $n$ buyers. Each buyer $i$ has a \textbf{private valuation} $\theta_i$ (willingness to pay). Utility is \textbf{quasi-linear}:
\[
  u_i = \begin{cases} \theta_i - p_i & \text{if } i \text{ wins and pays } p_i, \\ 0 & \text{otherwise.} \end{cases}
\]
Bidders do not know others' valuations but have beliefs described by a distribution $F$ (common knowledge). We assume no collusion (for now) and risk-neutral bidders.
\end{defbox}

\begin{center}
\renewcommand{\arraystretch}{1.3}
\begin{tabular}{llll}
\toprule
\textbf{Format} & \textbf{Type} & \textbf{Winner} & \textbf{Price paid} \\
\midrule
English         & Open, ascending  & Last remaining bidder & Own (final) bid \\
Dutch           & Open, descending & First to bid          & Own bid \\
First-price sealed-bid & Sealed & Highest bidder & Own bid \\
Vickrey (2nd-price)    & Sealed & Highest bidder & Second-highest bid \\
\bottomrule
\end{tabular}
\end{center}

\begin{tipbox}[Strategic Equivalences]
\begin{itemize}
  \item \textbf{Dutch $\equiv$ First-price sealed-bid:} same strategy space and BNE (both: trade off winning probability vs.\ payment). Dutch is sometimes called the ``open first-price auction.''
  \item \textbf{English $\approx$ Vickrey:} both result in the highest-valued bidder winning and paying $\approx$ the second-highest valuation. English is sometimes called the ``open second-price auction.''
\end{itemize}
\end{tipbox}

\separator

% --- Bayesian games ---
\subsubsection{Bayesian Games and Bayes-Nash Equilibrium}

\begin{defbox}[Bayesian Game]
A tuple $(N, (A_i, \Theta_i, u_i)_{i \in N}, \Pr)$ where:
\begin{itemize}
  \item $N$: set of agents.
  \item $A_i$: action set of player $i$.
  \item $\Theta_i$: \textbf{type space} of player $i$; player $i$ knows $\theta_i \in \Theta_i$ but not $\theta_{-i}$.
  \item $\Pr: \Theta_1 \times \cdots \times \Theta_n \to [0,1]$: \textbf{common prior} over type profiles.
  \item $u_i: (\times_{j} A_j) \times \Theta_i \to \mathbb{R}$: utility function.
\end{itemize}
A \textbf{strategy} is a mapping $s_i: \Theta_i \to A_i$.
\end{defbox}

\begin{defbox}[Bayes-Nash Equilibrium (BNE)]
A strategy profile $s = (s_1, \ldots, s_n)$ is a \textbf{BNE} if for every $i$, every $\theta_i$, and every $a_i \in A_i$:
\[
  \sum_{\theta_{-i}} \Pr(\theta_{-i})\, u_i(\theta_i, s_i(\theta_i), s_{-i}(\theta_{-i})) \;\geq\; \sum_{\theta_{-i}} \Pr(\theta_{-i})\, u_i(\theta_i, a_i, s_{-i}(\theta_{-i})).
\]
No player can improve expected utility by unilateral deviation, given beliefs about others' types.
\end{defbox}

\separator

% --- First-price BNE ---
\subsubsection{First-Price Sealed-Bid Auction: BNE}

\begin{thmbox}[BNE for First-Price Auction (Uniform Prior)]
Assume $n$ bidders with types $\theta_i \overset{\text{i.i.d.}}{\sim} U[0,1]$. The symmetric BNE strategy is:
\[
  s_i(\theta_i) = \frac{n-1}{n}\,\theta_i.
\]
\end{thmbox}

Key properties:
\begin{itemize}
  \item Bidders \textbf{shade their bid} below their true value. Shading decreases as $n \to \infty$.
  \item For $n = 2$: bid $\frac{1}{2}\theta_i$. For $n = 10$: bid $\frac{9}{10}\theta_i$.
  \item Strategy is \textbf{monotonically increasing} $\Rightarrow$ highest-valued bidder wins $\Rightarrow$ auction is \textbf{efficient}.
\end{itemize}

\begin{exbox}[First-Price Auction Payoff Matrix]
Two players: $\theta_1 = 4$, $\theta_2 = 2$. Integer bids $b_i \in \{0,1,2,3,4\}$, ties broken randomly.

\begin{center}
\small
\begin{tabular}{c|ccccc}
$\theta_1{=}4 \backslash \theta_2{=}2$ & $b_2{=}0$ & $b_2{=}1$ & $b_2{=}2$ & $b_2{=}3$ & $b_2{=}4$ \\
\hline
$b_1{=}0$ & 2, 1   & 0, 1     & 0, 0     & 0, $-1$  & 0, $-2$ \\
$b_1{=}1$ & 3, 0   & 1.5, 0.5 & 0, 0     & 0, $-1$  & 0, $-2$ \\
$b_1{=}2$ & 2, 0   & 2, 0     & 1, 0     & 0, $-1$  & 0, $-2$ \\
$b_1{=}3$ & 1, 0   & 1, 0     & 1, 0     & 0.5, $-0.5$ & 0, $-2$ \\
$b_1{=}4$ & 0, 0   & 0, 0     & 0, 0     & 0, 0     & 0, $-1$ \\
\end{tabular}
\end{center}

Note: this matrix assumes \emph{known} types. In practice, types are private $\Rightarrow$ need Bayesian game model.
\end{exbox}

\separator

% --- Vickrey auction ---
\subsubsection{Vickrey (Second-Price) Auction: Strategyproofness}

\begin{thmbox}[Truthful Bidding is Weakly Dominant in Vickrey Auction]
In the second-price sealed-bid auction, bidding $s_i(\theta_i) = \theta_i$ is a \textbf{weakly dominant strategy} for every bidder.
\end{thmbox}

\textbf{Proof.} Suppose bidder $i$ bids truthfully ($b_i = \theta_i$). Let $p = \max_{j \neq i} b_j$ (second-highest bid).

\emph{Case 1: $i$ wins ($\theta_i > p$).} Utility $= \theta_i - p > 0$.
\begin{itemize}
  \item Bidding \emph{lower}: same payment $p$ if still winning, but risks losing (utility drops to 0).
  \item Bidding \emph{higher}: no change (still wins at price $p$).
\end{itemize}

\emph{Case 2: $i$ loses ($\theta_i < p$).} Utility $= 0$.
\begin{itemize}
  \item Bidding \emph{lower}: no change (still loses).
  \item Bidding \emph{higher}: might win, but then pays $p > \theta_i$ $\Rightarrow$ negative utility.
\end{itemize}

In both cases, no deviation improves utility $\Rightarrow$ truthful bidding is weakly dominant. \qed

\begin{tipbox}[Why Vickrey Works]
The key insight: the \textbf{payment does not depend on the winner's bid}. Since $i$'s bid affects only whether they win (not what they pay), bidding truthfully maximises the chance of winning exactly when winning is profitable ($\theta_i > p$).
\end{tipbox}

\separator

% --- Revenue Equivalence ---
\subsubsection{Revenue Equivalence Theorem}

\begin{thmbox}[Revenue Equivalence Theorem]
Suppose valuations are drawn i.i.d.\ from a continuous distribution on $[L, H]$ and agents are risk-neutral. Then any two auctions that:
\begin{enumerate}
  \item In equilibrium, always allocate the item to the \textbf{highest-valued} bidder, and
  \item give a bidder with valuation $L$ an expected utility of $0$,
\end{enumerate}
yield the same \textbf{expected revenue} to the auctioneer.
\end{thmbox}

\textbf{Implication:} All four standard auctions (English, Dutch, first-price, Vickrey) are \textbf{revenue-equivalent} under these assumptions.

\separator

% --- Collusion ---
\subsubsection{Collusion}

\begin{defbox}[Collusion]
A group of bidders coordinates bids to reduce the price paid. The colluders split the surplus.
\end{defbox}

\textbf{Vulnerability by auction type:}
\begin{itemize}
  \item \textbf{English auction: vulnerable.} A \term{bidding ring} agrees not to bid against each other. Defection is visible, so the highest-valued ring member can retaliate by outbidding.
  \item \textbf{Dutch auction: resistant.} A ring member can secretly defect (others cannot respond before the item is sold).
\end{itemize}

\begin{exbox}[Collusion in Second-Price Auction (2 bidders only)]
Bidders 1, 2 with $\theta_1 > \theta_2$. Optimal collusion: bidder 2 bids 0, bidder 1 bids any $b_1 > 0$.

\textbf{Honest payoff:} $\theta_1 - \theta_2$ (bidder 1 wins, pays $\theta_2$). \textbf{Collusion payoff:} $\theta_1$ (bidder 1 wins, pays 0).

Extra surplus $= \theta_2$, split evenly: each gets $\frac{\theta_2}{2}$.
\end{exbox}

\begin{exbox}[Collusion with a Third Player]
Bidders 1, 2 collude ($\theta_1 > \theta_2$); bidder 3 draws $\theta_3 \sim U[0,1]$ independently.

Optimal strategy: $b_2 = 0$, $b_1 = \theta_1$ (play truthfully against bidder 3).

\textbf{Honest expected payoff} (bidder 1): $\displaystyle\int_0^{\theta_2}(\theta_1 - \theta_2)\,d\theta_3 + \int_{\theta_2}^{\theta_1}(\theta_1 - \theta_3)\,d\theta_3 = \frac{\theta_1^2 - \theta_2^2}{2}$.

\textbf{Collusion expected payoff} (bidder 1): $\displaystyle\int_0^{\theta_1}(\theta_1 - \theta_3)\,d\theta_3 = \frac{\theta_1^2}{2}$.

Extra surplus $= \frac{\theta_2^2}{2}$; each colluder gets $\frac{\theta_2^2}{4}$.
\end{exbox}

\separator

% --- Winner's curse ---
\subsubsection{Winner's Curse}

\begin{defbox}[Winner's Curse]
In \textbf{common-value auctions} (all bidders value the item equally, but estimate it with uncertainty), the winner is the bidder with the \textbf{highest estimate}---likely an \textbf{overestimate}. Formally:
\[
  \mathbb{E}[\text{value} - \text{price} \mid \text{win}] < 0.
\]
The curse is more severe with \textbf{more bidders} (the $n$-th order statistic increases with $n$).
\end{defbox}

\textbf{Mitigation:} Savvy bidders \textbf{shade their bids} downward, conditioning on the event of winning. The Revenue Equivalence Theorem implies rational bidders need never actually overpay.

\textbf{Real-world examples:} offshore oil field auctions (1950s--70s), spectrum auctions, IPOs, pay-per-click advertising.

\separator

% --- Cheating auctioneers ---
\subsubsection{Cheating Auctioneers}

\begin{itemize}
  \item \textbf{Vickrey auction:} auctioneer can lie about the second-highest bid (it is unverifiable).
  \item \textbf{English auction:} auctioneer can use \textbf{shill bidders} (fake bids to inflate the price).
\end{itemize}

\separator

% --- Multi-unit auctions ---
\subsubsection{Multi-Unit and Combinatorial Auctions}

\begin{defbox}[Multi-Unit Auction]
$k$ \textbf{identical} copies of a good. Payment rules:
\begin{itemize}
  \item \textbf{Discriminatory pricing (pay-your-bid):} each winner pays their own bid.
  \item \textbf{Uniform pricing:} all winners pay the same price (highest losing bid or lowest winning bid).
\end{itemize}
Bids can be \textbf{all-or-nothing} (fixed quantity) or \textbf{divisible} (accept fewer units at same per-unit price).
\end{defbox}

\begin{defbox}[Single-Unit Demand: $(k+1)$-st Price Auction]
If each bidder wants at most 1 unit: sell to the $k$ highest bidders, all paying the $(k+1)$-st highest bid. This is a special case of the \textbf{VCG mechanism} and is \textbf{efficient} and \textbf{dominant-strategy truthful}.
\end{defbox}

\begin{defbox}[Combinatorial Auction]
Multiple \textbf{heterogeneous} goods. Bidders' valuations depend on which \emph{subset} of goods they receive (complementarities and substitutes). Examples: spectrum auctions, energy markets, procurement.
\end{defbox}

\separator

% --- VCG ---
\subsubsection{Vickrey-Clarke-Groves (VCG) Mechanism}

\begin{defbox}[VCG Mechanism]
Given $n$ bidders, $m$ goods, type space $\Theta$, valuation $v: 2^G \times \Theta \to \mathbb{R}$, and feasible allocations $X \subseteq (2^G)^n$:

\textbf{Allocation rule} (efficient):
\[
  f(\hat{\theta}_1, \ldots, \hat{\theta}_n) \in \arg\max_{x \in X} \sum_{i \in N} v(x_i, \hat{\theta}_i).
\]

\textbf{Transfer rule} (externality payment):
\[
  t_i(\hat{\theta}) = \underbrace{\sum_{j \neq i} v\bigl(f(\hat{\theta})_j,\, \hat{\theta}_j\bigr)}_{\text{others' value when } i \text{ present}} - \underbrace{\max_{x \in X} \sum_{j \neq i} v(x_j, \hat{\theta}_j)}_{\text{others' best value without } i}.
\]

Agent $i$'s \textbf{payment} is $-t_i$: the \textbf{harm} $i$'s presence causes to other agents.
\end{defbox}

\begin{thmbox}[VCG Properties]
The VCG mechanism is:
\begin{enumerate}
  \item \textbf{Efficient:} maximises total welfare.
  \item \textbf{Dominant-strategy truthful:} reporting $\hat{\theta}_i = \theta_i$ is weakly dominant.
  \item \textbf{Individually rational:} $u_i \geq 0$ for all $i$.
  \item \textbf{Weakly budget-balanced:} total payments $\geq 0$ (mechanism never runs a deficit). Note: VCG is \emph{not} strongly budget-balanced (may have surplus).
\end{enumerate}
\end{thmbox}

\begin{tipbox}[VCG Reduces to Known Auctions]
\begin{itemize}
  \item \textbf{Single item:} VCG $=$ Vickrey (second-price) auction.
  \item \textbf{$k$ identical items, unit demand:} VCG $=$ $(k+1)$-st price auction.
\end{itemize}
\end{tipbox}

\separator

% --- Double auctions ---
\subsubsection{Continuous Double Auctions (CDA)}

\begin{defbox}[Continuous Double Auction]
Multiple buyers and sellers trade units of the same good. Each buyer has a \textbf{limit price} (max willing to pay); each seller has a \textbf{limit price} (min willing to accept). Agents bid at their own pace.

The \textbf{equilibrium price} is found where the \textbf{demand curve} (buyers' limit prices, sorted descending) crosses the \textbf{supply curve} (sellers' limit prices, sorted ascending). Trades occur between buyers and sellers whose limit prices are compatible at the equilibrium price.
\end{defbox}

\separator

% --- Worked examples ---
\subsubsection{Worked Examples from Exercises}

\begin{exbox}[VCG for 5 Identical Items, 7 Buyers (Exercise 1)]
\textbf{Setting:} 5 identical items, 7 buyers (unit demand). Valuations: $v_1 = 70$, $v_2 = 30$, $v_3 = 27$, $v_4 = 25$, $v_5 = 12$, $v_6 = 5$, $v_7 = 2$.

\textbf{(a) Assignment:} Efficient allocation gives items to the 5 highest-valued buyers: \textbf{buyers 1--5}.

\textbf{(b) Transfers/payments:} This is a $(k{+}1)$-st price auction. All winners pay the 6th-highest bid: $\mathbf{p = v_6 = 5}$.

Alternatively, computing VCG transfers directly: for each winner $i$, the payment equals the externality---the value the 6th bidder would have received without $i$'s presence. Since each winner displaces buyer 6 (the highest loser), each winner pays $v_6 = 5$.

Payments: buyers 1--5 each pay \textbf{5}. Total revenue = $5 \times 5 = 25$.

\textbf{(c) General rule:} With $k$ identical items and $n$ bidders (unit demand), VCG allocates to the $k$ highest bidders, each paying the $(k{+}1)$-st valuation (the highest losing bid). This is dominant-strategy truthful and efficient.
\end{exbox}

\begin{exbox}[VCG for Two Heterogeneous Items (Exercise 2)]
\textbf{Setting:} Items $\{A, B\}$, buyers $\{1, 2, 3\}$. Valuations:

\begin{center}
\begin{tabular}{lcccc}
\toprule
       & $\emptyset$ & $\{A\}$ & $\{B\}$ & $\{A,B\}$ \\
\midrule
Buyer 1 & 0 & 5  & 3  & 23 \\
Buyer 2 & 0 & 3  & 12 & 15 \\
Buyer 3 & 0 & 15 & 5  & 18 \\
\bottomrule
\end{tabular}
\end{center}

\textbf{(a) Efficient allocation:} Enumerate assignments and find the one maximising total value:
\begin{itemize}
  \item Give $\{A,B\}$ to buyer 1: total $= 23 + 0 + 0 = 23$.
  \item Give $A$ to buyer 3, $B$ to buyer 2: total $= 0 + 12 + 15 = 27$.
  \item Give $\{A,B\}$ to buyer 3: total $= 0 + 0 + 18 = 18$.
  \item Give $A$ to buyer 1, $B$ to buyer 2: total $= 5 + 12 + 0 = 17$.
  \item (Other combinations yield $\leq 27$.)
\end{itemize}

\textbf{Optimal:} $A \to$ buyer 3, $B \to$ buyer 2. Total welfare $= 27$.

\textbf{VCG payments:} For each buyer, payment $= -t_i = \text{(others' best without } i\text{)} - \text{(others' value with } i\text{)}$.

\emph{Buyer 2:} With buyer 2: others get $0 + 15 = 15$. Without buyer 2: best is $\{A,B\} \to$ buyer 1 giving $23$, or $A \to 3, B \to 1$ giving $3 + 15 = 18$, or $\{A,B\} \to 3$ giving $18$. Best without 2 $= 23$. Payment $= 23 - 15 = 8$.

\emph{Buyer 3:} With buyer 3: others get $0 + 12 = 12$. Without buyer 3: best is $\{A,B\} \to 1$ giving $23$, or $A \to 1, B \to 2$ giving $5 + 12 = 17$. Best without 3 $= 23$. Payment $= 23 - 12 = 11$.

\emph{Buyer 1:} Gets nothing in optimal allocation. Payment $= 0$.

\textbf{(b) Buyer 1's objection:} Buyer 1 values $\{A,B\}$ at 23, but gets nothing (total value of giving both to buyer 1 is only 23 vs.\ 27 for the optimal split). Despite having the highest individual bundle value, the efficient allocation serves the group better by splitting items.

\textbf{(c) Weakly budget-balanced?} Total payments $= 0 + 8 + 11 = 19 \geq 0$. Yes, the mechanism collects a surplus. \checkmark

\textbf{(d) Individually rational?} Buyer 1: $u_1 = 0 \geq 0$ \checkmark. Buyer 2: $u_2 = 12 - 8 = 4 \geq 0$ \checkmark. Buyer 3: $u_3 = 15 - 11 = 4 \geq 0$ \checkmark. Yes. \checkmark
\end{exbox}

\begin{exbox}[Continuous Double Auction (Exercise 3)]
\textbf{Setting:} 5 buyers with limit prices \pounds 2, \pounds 4, \pounds 6, \pounds 8, \pounds 10. Three sellers with limit prices \pounds 4, \pounds 8, \pounds 12.

Sort buyers descending: \pounds 10, \pounds 8, \pounds 6, \pounds 4, \pounds 2. Sort sellers ascending: \pounds 4, \pounds 8, \pounds 12.

Match where buyer limit $\geq$ seller limit:
\begin{itemize}
  \item Buyer (\pounds 10) -- Seller (\pounds 4): compatible.
  \item Buyer (\pounds 8) -- Seller (\pounds 8): compatible (equal).
  \item Buyer (\pounds 6) -- Seller (\pounds 12): \textbf{not compatible}.
\end{itemize}

Demand curve crosses supply curve between the 2nd and 3rd trades. The \textbf{equilibrium price} is \pounds\textbf{8} (where the last feasible trade occurs).

\textbf{Transactions:} Two trades at \pounds 8---the buyer with limit \pounds 10 and buyer with limit \pounds 8 trade with the sellers with limits \pounds 4 and \pounds 8.
\end{exbox}

\separator

% --- Summary table ---
\subsubsection{Summary: Auction Properties}

\begin{center}
\renewcommand{\arraystretch}{1.3}
\begin{tabular}{lcccc}
\toprule
\textbf{Property} & \textbf{First-price} & \textbf{Vickrey} & \textbf{English} & \textbf{Dutch} \\
\midrule
Dominant-strategy truthful & $\times$ & \checkmark & $\approx$ \checkmark & $\times$ \\
Efficient (in eq.)         & \checkmark & \checkmark & \checkmark & \checkmark \\
Revenue equivalent         & \checkmark & \checkmark & \checkmark & \checkmark \\
Resistant to collusion     & Partial & $\times$ & $\times$ & \checkmark \\
Resistant to cheating auctioneer & \checkmark & $\times$ & $\times$ & \checkmark \\
\bottomrule
\end{tabular}
\end{center}

\begin{center}
\renewcommand{\arraystretch}{1.3}
\begin{tabular}{lccccc}
\toprule
\textbf{Mechanism} & \textbf{Efficient} & \textbf{Truthful} & \textbf{IR} & \textbf{W.~Budget-Bal.} & \textbf{S.~Budget-Bal.} \\
\midrule
VCG         & \checkmark & \checkmark & \checkmark & \checkmark & $\times$ \\
$(k{+}1)$-st price & \checkmark & \checkmark & \checkmark & \checkmark & $\times$ \\
Vickrey     & \checkmark & \checkmark & \checkmark & \checkmark & $\times$ \\
\bottomrule
\end{tabular}
\end{center}
